% ============================================================
%  第三卷 第一章 量子力学的基本概念（§1–§7）
%  底本：高教社中文版第6版（严肃 译，喀兴林 校）
%  本章无插图、无习题（与原书一致）；正文圈码 18 处已全部转录为
%  \footnote（脚注代理，2026-10-11，按原书页面 p014–p034 逐页核录：
%  §1 六处 p014×2/p015×3/p016，§2 三处 p018×2/p019，§3 四处
%  p022×2/p023×2，§5 三处 p029/p030×2，§6 一处 p031，§7 一处 p034；
%  脚注内公式照原书排入，含 §3 脚注②线性算符条件两式——该两式
%  曾被 OCR 误排正文（行551）而由排版代理删除，现已归位脚注）。
%  已修正的 OCR/原书错误（均对照 原书转换/pages/p-014…p-034.png 核实）：
%   1. 波函数记号：OCR ψ/Ψ 混排且多处误作粗体 \pmb{\psi}；按原书一律用
%      大写 \varPsi（态波函数与本征函数 Ψₙ 均大写）；q、f、a_n 等为斜体，
%      (2.3)(2.4) 的 q_1,q_2 及 (3.10) 的 K(q,q′) OCR 误作 \boldsymbol。
%   2. |Ψ|^2 类乱码：OCR "141^2"、"1Ψ1^2"、"|Ψ12"、"Ψ_f1^2" 等一律还原
%      为 |\varPsi|^2（§2 多处）。
%   3. 行450 乱句按原书 p.7 重建："最后，当系统处于波函数 Ψ=Ψₙ 的态中时，
%      fₙ 值的概率必须等于1，当波函数 Ψ 的展开式(3.2)中不含 Ψₙ 时……"。
%   4. 行551 线性算符条件 f̂(Ψ₁+Ψ₂)=f̂Ψ₁+f̂Ψ₂、f̂(aΨ)=af̂Ψ 系原书脚注②
%      文字窜入正文（行510 圈码②所指），正文删除。
%   5. (3.7) OCR \tilde{f}→\bar{f}；(3.16) OCR \tilde{\hat f}^{·}→
%      \tilde{\hat f}^{*}。
%   6. 行504/510/562/574 的 \textcircled{} 为脚注圈码①②，原样保留。
%   7. §3 正交性证明：OCR "fn和f" 补正为 fₙ 和 fₘ；ΨₙΨₙ*→ΨₙΨₘ* 等
%      三处下标 m/n 混淆按原书改正。
%   8. (4.4) OCR 乱码按原书重建为 \widetilde{\hat{f}\hat{g}}
%      =\tilde{\hat{g}}\tilde{\hat{f}}；行675/681 积分式内 \hat{\psi}、
%      \hat{\phi} 的帽子系 OCR 噪声，按原书为 Ψ、Φ；行681 中项放大核对
%      确认为 (f̃̂Ψ)(ĝΦ)。
%   9. 行696 "(f̂ĝ)^*=ĝf̂"→(f̂ĝ)⁺=ĝf̂（原书为厄米共轭上标 +）。
%  10. §6 普朗克常量：OCR 正文误作 h，按原书一律为 ħ（\hbar）；行920
%      夹杂乱码 "$\times\varphi$" 删。
%  11. §7 行944 "\pmb{\theta}_n"→g_n、"读数"g_n 引号错位按原书改；
%      行1006 Ψₚ→Ψₙ；行1010 积分内 ψₘ*→Ψₘ*（大写）。
%  12. 圈码噪声：行414 "第一个⃝1子系统"、行1010 "算出其⃝1概率" 处的
%      圈码经原书页面核对为窜入噪声，删除；正文实有圈码 18 处未丢一个。
%  13. 半角标点 , ; : ( ) 按文体约定改全角；数学环境内保持 ASCII。
%  遗留问题：
%   - 行344 圈码①在"1927年"之后，原书如此（脚注①讲海森伯 1925—1926
%     年建立完整数学表述），非 OCR 错位；
%   - (3.16) 式尾标点原书版面不清，按 OCR 不加。
% ============================================================
\chapter{量子力学的基本概念}

\section{不确定性原理}

每当我们试用经典力学和经典电动力学阐释原子现象时，总会得出与实验有明显矛盾的结论.最明显的例子，是把通常的电动力学用于电子绕原子核作经典轨道运动的原子模型.当电子作这种运动的时候，它和任何带电粒子的加速运动一样，会不断地辐射电磁波.由于这种辐射，电子便会丧失能量，这将使它最终落入原子核中.故按经典电动力学看来，原子是不稳定的，但这与事实完全不符.

理论与实验之间如此深刻的矛盾，表明有必要建立一种适用于原子现象（即质量极小的一些粒子在极短间距内所发生的现象）的理论，需要根本改变基本的物理概念和定律.

为方便计，我们拿实验上观察到的电子衍射现象\footnote{电子衍射现象实际上是在量子力学建立以后才发现的.但在我们的论述中，我们不去拘泥于理论的历史发展顺序，而是尽量采用这样的讲法，使得量子力学的基本原理与实验现象之间的联系表达得最为清楚.}，作为阐明这种根本改变的出发点.当一均匀电子束穿过一块晶体时，发现出射波呈现一种强弱交替的图样，完全类似于电磁波衍射中所观察到的衍射图样，由此可见，在一定的条件下，粒子（此例中为电子）的行为中会表现出属于波动过程的特征.

这种现象与习常的运动观念之间的矛盾，究竟尖锐到什么地步，最好用以下的假想实验说明，它是晶体的电子衍射实验的一种抽象化.设想有一块电子不能穿透的屏板，板上开有两道狭缝.让电子束\footnote{假定粒子束是如此稀疏，以致粒子间的相互作用可以略去不计.}通过其中的一个狭缝（遮住另一个），则在狭缝后放置的连续幕上可以得到某一强度分布图样；应用同样的方法，遮闭第一个狭缝并打开第二个，可得另一个图样.现在让电子束同时通过这两个狭缝，我们根据通常的经典观念，一定会设想所得的图样不过是原先两个图样的简单叠合：因为每一个电子都沿自己的轨道运动，只通过狭缝之一，而不会影响正在通过另一个狭缝的电子.可是，电子衍射现象表明，由于干涉作用，我们所得的衍射图样实际上并不等于每一个单狭缝所分别给出的那两个衍射图样之和.十分明显，这个结果无法与电子的轨道运动观念相协调.

因此，统辖原子现象的力学——量子力学或波动力学——必须建立在与经典力学根本不同的运动观念的基础之上.量子力学中并不存在粒子轨道之类的概念.这就构成了1927年\footnote{值得指出，量子力学的完整数学表述，是在不确定性原理发现之前，由W.海森伯和E.薛定谔在1925—1926年间建立起来的.不确定性原理体现了这一数学表述的物理内容.}W.海森伯所发现的量子力学基本原理之一所谓不确定性原理的主要内容.

从抛弃经典力学的习常观念这一角度来讲，不确定性原理的内容也许可以说是消极的.诚然，这个原理本身，还不足以作为建立新的粒子力学的基础.这样一种新的理论，自应建立在若干积极论断的基础上，这将在以后（§2）讨论.但是，为了能够表述这些论断起见，我们有必要首先弄清量子力学所面临的问题的提法.为此，我们先来考察一下量子力学和经典力学内在关系间的特殊性质.

凡是一个更为普遍的理论，往往可用完整的逻辑形式表述出来，并且独立于那些作为它的极限情形的较窄理论.例如相对论力学可以建立在自己的基本原理的基础上，无需参考牛顿力学.可是，当我们表述量子力学的基本概念时，原则上却不能不用到经典力学.一个电子\footnote{为简便计，在本节及以后各节中，凡是讲到“一个电子”的地方，可以一般地理解为一个具有量子特性的任何客体，即指不服从经典力学而服从量子力学的粒子或粒子系统.}没有确定的轨道，这一事实本身意味着这个电子也不会有其它什么动力学标志\footnote{我们所指的是标志电子运动的那些量，而不是指标志粒子本身的电荷、质量等等参量.}.于是就很清楚，对于一个只包含量子客体的系统来讲，势必完全不可能建立起任何逻辑上独立的力学.对电子运动作出定量描述的可能性，要求同时存在一些物理客体，这些物理客体在足够精确的范围内服从经典力学.如果一个电子和这样的“经典客体”相作用，后者的状态一般讲来会有所变化.这一变化的性质及大小依赖于电子的状态，从而就可以用来定量描述电子的状态.

因而，“经典客体”通常称为仪器，它和电子的作用就称为测量.但是有必要强调指出，我们在这里根本没有讨论物理观测者所参与的“测量”过程.量子力学中所谓的测量，我们总是把它理解为与任何观测者无关的发生于经典客体和量子客体之间的任一相互作用过程.测量概念在量子力学中的重要性是由N.玻尔所阐明的.

我们已把“仪器”定义为在足够精确范围内服从经典力学的一个物理客体.例如一个质量足够大的物体.但不能因此认为仪器必然是宏观的.在一定条件下，微观客体也能起部分仪器的作用，因为“具有足够的精确度”这一概念取决于所涉及的具体问题.例如威耳逊云室中的电子运动，可用它所遗留的云迹来观察，这种云迹的粗细远大于原子尺度；当用这样低的精确度确定轨道时，电子完全是一个经典客体.

由此可见，量子力学在物理理论中占有一个很不平常的地位；它把经典力学作为一种极限情形而包含在内，但在它的自身表述中，同时又需要这一极限情形.

现在可以来表述一下量子力学问题的提法.一种典型的问题是：用前次测量的已知结果，去预断下一次的测量结果.除此以外，我们以后将看到，量子力学中的各种物理量（例如能量）所能采取的数值，即作为该量的测量结果所能得到的数值，它们的值域和经典力学相比一般讲来是受限制的.量子力学方法必须告诉我们怎样来确定它们的各种允许值.

量子力学中的测量过程具有一种十分重要的特性：它总是要影响到被测的电子，并且在给定的测量精确度范围内，原则上不可能使这种影响变得任意小.测量得愈精确，它所给予的影响就愈大，只有在精确度极低的测量中，被测客体所受的影响才能很小.测量的这种性质，逻辑上是由于电子的动力学标志仅仅作为测量本身的结果才能表现出来；十分明显，如果测量过程对客体的影响可以任意地小，这就意味着被测的那个量本身具有一个和测量无关的定值.

在各种测量中，电子的坐标测量具有基本的意义.在量子力学的适用限度内，对一个电子所施行的坐标测量\footnote{我们再强调一遍，所谓“施行测量”是指一个电子和一个经典“仪器”的相互作用，丝毫也没有预设外界观测者的存在.}总是可以达到需要的任何精确度.

现在假定对一个电子的坐标相继测量了许多次，每次相隔的时间固定为$\Delta t$.这些测量结果，一般讲来，并不位于一条光滑的曲线上.而是相反，测量得愈精确，这些结果会变化得愈不连续愈不规则；正好和电子不存在轨道的概念相一致.只有在极为粗略地测量电子坐标的情形下，例如，在威耳逊云室中根据蒸气凝成的液滴确定电子坐标的情形下，才会得到一条相当光滑的轨道.

现在让测量的精确度保持不变，我们把测量之间的时间间隔$\Delta t$加以缩短，那么，相邻的测量当然就会给出坐标的相邻值.一系列相继测量之后所得的结果，虽然都会落到某一很小的空间范围内，可是它们将在这个区域内毫无规则地分布着，并不位于任一光滑曲线上.特别是$\Delta t$趋向于零时，相继测量的结果完全不趋于同一直线上.

这种情况表明，量子力学中并不存在经典意义下的粒子速度概念，经典的粒子速度，就是指两个时刻的坐标之差除以这两个时刻的时间间隔$\Delta t$后当$\Delta t$趋向于零时所得的极限.不过，以后我们会看到，量子力学中可以建立一个合理的定义，用来表示某一给定时刻的粒子速度，并且当量子力学转向经典力学时，这个速度也随之转为经典的速度.可是，在经典力学中，一个粒子在任一给定时刻可以同时具有确定的坐标和确定的速度，而在量子力学中，情况则完全不同.如果测量结果发现电子具有确定的坐标，那么，它就不可能同时具有任何确定的速度.反之，具有确定速度的电子，就不可能具有确定的空间位置.事实上，坐标和速度的同时存在就意味着一条确定轨道的存在，这正是电子所没有的.由此可见，在量子力学中，一个电子的坐标和速度是两个不能同时确切测量的量，也就是两个不能同时具有定值的量.我们也可以说，电子的坐标和速度是两个无法同时存在的量.以后我们还要导出一个定量的关系式，用来判断坐标和速度同时进行非精确测量的可能性.

在经典力学中，物理系统的完全描述是通过某一时刻给定其中的全部坐标和速度来实现的.运动方程根据这些初始条件就可以完全确定系统此后所有时刻的行为.而在量子力学中这种描述在原则上是不可能的，因为坐标和与其相应的速度不可能同时确定.于是，量子系统的状态是用远比经典力学为少的数值描写的.这就是说，对量子系统的状态的描写不如经典系统那样详细.

由此得出关于量子力学所能作出的预言的性质的一个非常重要的结论：和经典力学的描述能够准确地预言系统此后的运动相比，量子力学对力学系统的描述是不够详细的，显然做不到像经典力学那样.这就是说，即使电子处于一个描述得最完全的态中，在以后的时刻它的行为在原则上也是不唯一的.因此量子力学对于电子此后的行为不可能给出确切的预言.处于给定起始状态的电子，此后对它进行测量时可以得出各种各样的结果.量子力学的任务只是给出测量得到这一结果或那一结果的概率.当然，在有的情况下测量得到某一结果的概率可能等于1，这时过渡到确定性，这一测量结果是单值的.

量子力学中所有的测量过程可以分成两类.其中有一类居多数，这类测量不管系统处于什么状态，都不能测得唯一的肯定结果.另外一类测量，它的每一种可能结果都能从某种相应状态中肯定地测得.后一类测量在量子力学中占有重要的地位，称为可以预断的测量.由这种测量所确定的状态，它的定量标志称为量子力学中的物理量.如果在某一态中，某种测量总是给出唯一的肯定结果，我们就说该态中相应的物理量具有定值.今后我们对“物理量”一词，总是理解为此处所指的含义.

今后我们一再会看到，量子力学中远非任意一组物理量都能同时测量，即都能同时具有定值.我们早已讲过一个例子，就是一个电子的坐标和速度.在量子力学中起着巨大作用的是具有下述性质的一组物理量：这组量能够同时测量，并且当它们同时具有定值的时候，再也没有别的物理量（只要不是这一组量的函数）能在该态中再具有定值，我们把这样的一组物理量称为一个完全集合.

电子状态的任何描述全都来自某种测量结果.现在来讲一下量子力学中对一个状态进行完全描述的含义.完全描述的态是由物理量的某一完全集合同时测量的结果所产生的.根据这样的测量结果，我们就能确定下一次任何测量中各种所得结果的概率，并与首次测量（完全测量）之前电子的历史无关.

今后（§14除外）我们总是把一个量子系统的状态理解成为这种完全描述的态.

\section{叠加原理}

量子力学中的运动概念和经典力学相比较有了根本性的改变，当然，这就要求理论的数学表述作出同样的根本改变.对此，我们必须首先考虑量子力学中态的描述方法.

我们用$q$表示量子系统坐标的集合，用$\mathrm{d}q$表示这组坐标的微分的乘积.$\mathrm{d}q$通常称为该系统位形空间中的一个体积元；对单粒子来讲，$\mathrm{d}q$等同于普通空间中的一个体积元$\mathrm{d}V$.

量子力学的数学表述基于这样的一个命题：在某一给定时刻，一个系统的状态可以用一个确定的坐标函数$\varPsi(q)$（通常为复函数）来描述.这个函数的模量平方确定了坐标值的概率分布：对系统进行坐标测量时，测量所得诸值处于位形空间的$\mathrm{d}q$体积元内的概率等于$|\varPsi|^2\,\mathrm{d}q$.$\varPsi$函数称为该系统的波函数\footnote{波函数是由薛定谔于1926年首先引入量子力学的.}.

知道波函数后，我们还能在原则上算出任何其它测量（不一定是坐标测量）结果的概率.所有这些概率都可由$\varPsi$和$\varPsi^{*}$的双线性表式所确定.这种表式的最一般形式为
\begin{equation}
  \iint \varPsi(q)\varPsi^{*}(q')\phi(q,q')\,\mathrm{d}q\,\mathrm{d}q',
\end{equation}
其中的$\phi(q,q')$函数依赖于测量的结果及性质，积分则遍及整个位形空间.坐标值的概率式$\varPsi\varPsi^{*}$本身，也是属于这种类型的一个表式\footnote{(2.1)式中当$\phi(q,q')=\delta(q-q_0)\delta(q'-q_0)$时可得$\varPsi(q_0)\varPsi^{*}(q_0)$，其中的$\delta$代表后面§5中定义的德耳塔函数；$q_0$为我们欲求其概率的一组坐标值.}.

一般讲来，系统的状态及其波函数会随时间变化.在这种意义下，波函数也可以看成是时间的函数.如果某一起始时刻的波函数是已知的，那么，根据状态的完全描述这一概念本身所具的含义，在原则上可确定此后每一时刻的波函数.波函数对时间的具体依赖关系，将由以后导出的方程式来确定.

根据定义，一个系统的各种可能坐标值的概率总和必须等于1.故$|\varPsi|^2$对整个位形空间的积分结果必须等于1：
\begin{equation}
  \int|\varPsi|^2\,\mathrm{d}q=1.
\end{equation}
这个等式称为波函数的归一化条件.如果$|\varPsi|^2$的积分是收敛的，那么，只要适当选择函数$\varPsi$前的常数系数，总能使$\varPsi$得到归一化.但是，我们在以后还会碰到$|\varPsi|^2$的积分为发散的波函数.以致$\varPsi$无法用条件 (2.2) 加以归一化.当然，这种情形下的$|\varPsi|^2$并不代表坐标的绝对概率值，但是，在位形空间中两个不同点处的$|\varPsi|^2$的比值，可给出这两处坐标值的相对概率.

凡用波函数算出并且具有直接物理意义的各种量，都呈 (2.1) 式的形式，式中的$\varPsi$总是跟$\varPsi^{*}$乘在一起，由于这一点，归一化的波函数显然可以包含一个具有$e^{i\alpha}$（$\alpha$为任一实数）形式的不确定的常数相因子，这个因子的模量等于1.这种不确定性原则上是无法消除的；但它无关紧要，因为它并不影响任何物理结果.

量子力学的积极内容是建立在有关波函数性质的一系列假定的基础之上的.这些假定如下：

设在波函数为$\varPsi_1(q)$的态中进行某种测量可以获得可靠的肯定结果（称为结果1），而在$\varPsi_2(q)$的态中进行这种测量也可以获得可靠的肯定结果2.那么可以假定，在$\varPsi_1$和$\varPsi_2$的任一线性叠加所给出的态中，即在任一具有$c_1\varPsi_1+c_2\varPsi_2$函数形式（其中$c_1$和$c_2$为常数）的态中，进行该种测量所得的结果或者是1，或者是2.此外，还可进一步假定，只要以上两个态的时间依赖关系是已知的，也就是一个由函数$\varPsi_1(q,t)$给出，另一个由函数$\varPsi_2(q,t)$给出，那么，它们的任一线性叠加也给出了这个叠加态的可能的时间依赖关系.以上这些假定，构成了量子力学的一个首要原理，称为状态叠加原理.从这个原理可以立刻知道，波函数所满足的一切方程必须对$\varPsi$保持线性.

现在让我们考虑一个系统，它由两个子系统所组成，并且假定这个系统处于这样的状态，它的每一个子系统都是完全描述的\footnote{当然，这意味着整个系统的态也是完全描述的.但应强调指出相反的情况并不成立：如果整个系统的态是完全描述的，一般来讲，它并不能由此确定各个个别部分的态（又见§14）.}.那么我们可以断言，第一个子系统的$q_1$坐标概率和第二个子系统的$q_2$坐标概率彼此无关，从而整个系统的概率分布必然等于这两个子系统的概率分布的乘积.这就是说，该系统的波函数$\varPsi_{12}(q_1,q_2)$可以表为这两个子系统波函数$\varPsi_1(q_1)$和$\varPsi_2(q_2)$的乘积：
\begin{equation}
  \varPsi_{12}(q_1,q_2)=\varPsi_1(q_1)\varPsi_2(q_2),
\end{equation}
如果这两个子系统没有相互作用，那么，整个系统及其两部分之间的上述波函数关系式在此后各时刻仍将保持不变，即可写成
\begin{equation}
  \varPsi_{12}(q_1,q_2,t)=\varPsi_1(q_1,t)\varPsi_2(q_2,t).
\end{equation}

\section{算符}

现在来考虑某一标志量子系统状态的物理量$f$.严格讲来，我们下面所讨论的量不是一个，而是同时有一组这种物理量的完全集合.但由于对问题的讨论没有多大差别，为简明计，下面只讲一个物理量.

在量子力学中，一个给定物理量所能取的数值称为它的本征值，这些数值的集合则称为该量的值谱.在经典力学中，一般讲来，物理量具有连续值.量子力学中也有一些物理量（例如坐标），它们的本征值具有连续的值域；这种情形下的本征值称为具有连续谱.可是，量子力学中除了这些量以外，还有其它一些量，它们的本征值是一组离散的数值；这种情形我们称为具有离散谱.

为简单计，我们假定所考虑的量$f$具有离散谱；连续谱的情形将在§5中讨论.$f$的本征值用$f_n$表示，下标$n$可取$0,1,2,3,\cdots$等值.我们还用$\varPsi_n$表示系统的一个状态波函数，该态中的$f$具有确定的$f_n$值.这个波函数$\varPsi_n$称为所给物理量$f$的本征函数.假定每一个这样的波函数已经归一化，即有
\begin{equation}
  \int|\varPsi_n|^2\,\mathrm{d}q=1.
\end{equation}
假定该系统处于某一波函数为$\varPsi$的任意状态中，对之进行物理量$f$的测量结果可得$f$的某些本征值$f_n$.根据状态叠加原理可以断言，波函数$\varPsi$必然是这样一些本征函数$\varPsi_n$的线性组合，这些$\varPsi_n$所对应的各个$f_n$值都能在$\varPsi$态中测到，且其概率不等于零.因此在一般情形下，任一态的$\varPsi$函数可表成下列级数形式：
\begin{equation}
  \varPsi=\sum_n a_n\varPsi_n,
\end{equation}
式中对所有的$n$求和，$a_n$是一些常系数.

由此得出结论，任何波函数可以用任一物理量的一套本征函数来展开.使上述展开式得以成立的那一套函数称为一个完备组（或封闭组）.

当系统处于波函数为$\varPsi$的态时，展开式 (3.2) 足以确定在该态中找到物理量$f$具有任一给定$f_n$值的概率（即测量结果为$f_n$值的概率）.正如上节所述，这些概率应由$\varPsi$和$\varPsi^{*}$的某种双线性表式所确定，这种表式对$a_n$和$a_n^{*}$而言因而也是双线性的.当然，这样的表式还必须是正量.最后，当系统处于波函数$\varPsi=\varPsi_n$的态中时，$f_n$值的概率必须等于1，当波函数$\varPsi$的展开式 (3.2) 中不含$\varPsi_n$时$f_n$值的概率必须等于零.能够满足上述诸条件的唯一正量只能是系数$a_n$的模量平方.我们就得到这样的结论，展开式 (3.2) 中每一个系数的模量平方值$|a_n|^2$，确定了波函数为$\varPsi$的态中物理量$f$具有相应值$f_n$的概率.各种$f_n$值的概率总和必须等于1；换句话说，以下关系式必须成立：
\begin{equation}
  \sum_n|a_n|^2=1.
\end{equation}
如果函数$\varPsi$未经归一化，(3.3) 式也就不再成立.此时$\sum_n|a_n|^2$将由$\varPsi$和$\varPsi^{*}$的某一双线性表式所确定，这个表式当$\varPsi$归一化后必须等于1.这样的表式只能是积分式$\int\varPsi\varPsi^{*}\,\mathrm{d}q$.故下式必须成立：
\begin{equation}
  \sum_n a_na_n^{*}=\int\varPsi\varPsi^{*}\,\mathrm{d}q.
\end{equation}
另一方面，如果我们用$\varPsi$去乘$\varPsi^{*}$（$\varPsi$的共轭复函数）的展开式$\varPsi^{*}=\sum_n a_n^{*}\varPsi_n^{*}$，积分之，可得
\begin{equation*}
  \int\varPsi\varPsi^{*}\,\mathrm{d}q=\sum_n a_n^{*}\int\varPsi_n^{*}\varPsi\,\mathrm{d}q,
\end{equation*}
将此式和 (3.4) 式比较，我们有
\begin{equation*}
  \sum_n a_na_n^{*}=\sum_n a_n^{*}\int\varPsi_n^{*}\varPsi\,\mathrm{d}q,
\end{equation*}
由此可以导出下列公式：
\begin{equation}
  a_n=\int\varPsi\varPsi_n^{*}\,\mathrm{d}q,
\end{equation}
这个公式确定了$\varPsi$函数对$\varPsi_n$本征函数组展开时的系数$a_n$.

如果把 (3.2) 式代入上式，可得
\begin{equation*}
  a_n=\sum_m a_m\int\varPsi_m\varPsi_n^{*}\,\mathrm{d}q,
\end{equation*}
由上式显然可知，本征函数组必须满足以下条件：
\begin{equation}
  \int\varPsi_m\varPsi_n^{*}\,\mathrm{d}q=\delta_{nm},
\end{equation}
$n=m$时其中的$\delta_{nm}=1$，$n\neq m$时$\delta_{nm}=0$.$m\neq n$时乘积$\varPsi_m\varPsi_n^{*}$的积分等于零，这一事实称为$\varPsi_n$函数组的正交性.因此，$\varPsi_n$本征函数组构成一套正交的归一化的完备函数组（或简称为正交归一系）.

现在来引入物理量$f$在某一给定态中的平均值$\bar{f}$的概念.根据通常的平均值定义，我们把平均值$\bar{f}$定义为该量的所有本征值$f_n$分别乘以相应的概率值$|a_n|^2$后相加所得的总和，即
\begin{equation}
  \bar{f}=\sum_n f_n|a_n|^2.
\end{equation}
我们来写出$\bar{f}$的一个表式，这个表式中不含$\varPsi$的展开系数$a_n$而只含$\varPsi$函数本身.鉴于 (3.7) 式中出现乘积$a_na_n^{*}$，显然，欲求的表式对$\varPsi$和$\varPsi^{*}$必须是双线性的.我们引入某种数学算符，用$\hat{f}$表示之\footnote{我们约定，凡在字母上加一符号“\textasciicircum{}”的都代表算符.}，并定义如下.令$(\hat{f}\varPsi)$为算符$\hat{f}$作用于函数$\varPsi$后所得的结果.我们定义的$\hat{f}$，是使$(\hat{f}\varPsi)$和共轭复函数$\varPsi^{*}$相乘后的积分结果等于平均值$\bar{f}$：
\begin{equation}
  \bar{f}=\int\varPsi^{*}(\hat{f}\varPsi)\,\mathrm{d}q.
\end{equation}
很易证明，$\hat{f}$在一般情形下是一个线性\footnote{一个算符具有以下性质时，称为线性的：
\begin{equation*}
  \hat{f}(\varPsi_1+\varPsi_2)=\hat{f}\varPsi_1+\hat{f}\varPsi_2
  \ \text{及}\ \hat{f}(a\varPsi)=a\hat{f}\varPsi,
\end{equation*}
其中$\varPsi_1$和$\varPsi_2$为任意函数，$a$为任意常数.}积分算符.实际上，应用$a_n$的表式 (3.5)，可以把 (3.7) 式的平均值定义改写成
\begin{equation*}
  \bar{f}=\sum_n f_na_na_n^{*}
  =\int\varPsi^{*}\left(\sum_n a_nf_n\varPsi_n\right)\mathrm{d}q,
\end{equation*}
和 (3.8) 式比较，可以看出算符$\hat{f}$作用于函数$\varPsi$的结果呈以下形式：
\begin{equation}
  (\hat{f}\varPsi)=\sum_n a_nf_n\varPsi_n,
\end{equation}
如果把 (3.5) 式的$a_n$代入上式，我们就发现$\hat{f}$是一个以下形式的积分算符：
\begin{equation}
  (\hat{f}\varPsi)=\int K(q,q')\varPsi(q')\,\mathrm{d}q',
\end{equation}
其中的函数$K(q,q')$（称为该算符的核）为
\begin{equation}
  K(q,q')=\sum_n f_n\varPsi_n^{*}(q')\varPsi_n(q).
\end{equation}
因此，量子力学中的每一个物理量都有一个确定的线性算符与之相应.

由 (3.9) 式可知，如果函数$\varPsi$就是本征函数$\varPsi_n$之一（此时所有的$a_n$除了一个以外都等于零），那么，用算符$\hat{f}$作用之后，就等于这个函数乘上它的相应本征值$f_n$：
\begin{equation}
  \hat{f}\varPsi_n=f_n\varPsi_n
\end{equation}
[此后凡不会发生混淆之处，常把$(\hat{f}\varPsi)$表式中的括号略去；将该算符视为作用在它后面所跟的表式上].因此我们可以说，所给物理量$f$的诸本征函数均为下列方程的解：
\begin{equation*}
  \hat{f}\varPsi=f\varPsi,
\end{equation*}
其中的$f$是一个常数，当上述方程具有满足所需条件的解答时，这个常数所能采取的诸数值即为本征值.以后将会看到，许多物理量的算符形状可以通过直接的物理考虑确定出来，上面所讲的算符性质，使我们有可能通过解$\hat{f}\varPsi=f\varPsi$方程而求出其本征函数及本征值.

一个实物理量的本征值及其在每个态中的平均值都是实数.这一点限制了它的相应算符.令 (3.8) 式等于它的复共轭式，可得以下关系：
\begin{equation}
  \int\varPsi^{*}(\hat{f}\varPsi)\,\mathrm{d}q
  =\int\varPsi(\hat{f}^{*}\varPsi^{*})\,\mathrm{d}q,
\end{equation}
其中的$\hat{f}^{*}$代表$\hat{f}$的复共轭算符\footnote{对算符$\hat{f}$，如果有$\hat{f}\varPsi=\varPhi$，则复共轭算符$\hat{f}^{*}$可通过$\hat{f}^{*}\varPsi^{*}=\varPhi^{*}$加以定义.}.一般讲来，这个式子对任意的线性算符不一定成立，因而它对算符$\hat{f}$的形式是一种限制.对任意的算符$\hat{f}$讲来，我们可以求出它的转置算符$\tilde{\hat{f}}$，其定义如下：
\begin{equation}
  \int\varPhi(\hat{f}\varPsi)\,\mathrm{d}q
  =\int\varPsi(\tilde{\hat{f}}\varPhi)\,\mathrm{d}q,
\end{equation}
式中的$\varPsi$和$\varPhi$是两个不同的函数.如果令$\varPhi$等于$\varPsi$的共轭复函数$\varPsi^{*}$，再和 (3.13) 式比较，我们就有
\begin{equation}
  \tilde{\hat{f}}=\hat{f}^{*},
\end{equation}
满足这个条件的算符称为厄米算符\footnote{对(3.10)形式的线性积分算符而言，厄米条件相当于该算符的核必须满足$K(q,q')=K^{*}(q',q)$.}.这样一来，在量子力学的数学表述中，和实物理量相对应的算符就必须是厄米算符.

我们还可以纯形式地考虑复物理量，即其本征值为复数的物理量.假定$f$是这样一个量.则可引进其共轭复量$f^{*}$，它的本征值和$f$的本征值成共轭复数关系.我们用$\hat{f}^{+}$代表$f^{*}$所对应的算符.$\hat{f}^{+}$称为算符$\hat{f}$的厄米共轭算符，一般讲来它不同于复共轭算符$\hat{f}^{*}$：$f^{*}$在$\varPsi$态中的平均值为
\begin{equation*}
  \bar{f}^{*}=\int\varPsi^{*}\hat{f}^{*}\varPsi\,\mathrm{d}q.
\end{equation*}
另一方面，我们有
\begin{equation*}
  (\bar{f})^{*}=\left[\int\varPsi^{*}\hat{f}\varPsi\,\mathrm{d}q\right]^{*}
  =\int\varPsi\hat{f}^{*}\varPsi^{*}\,\mathrm{d}q
  =\int\varPsi^{*}\tilde{\hat{f}}^{*}\varPsi\,\mathrm{d}q
\end{equation*}
以上两式相等，得
\begin{equation}
  \hat{f}^{+}=\tilde{\hat{f}}^{*}
\end{equation}
显然，$\hat{f}^{+}$一般讲来并不等于$\hat{f}^{*}$.

条件 (3.15) 现在可写成
\begin{equation}
  \hat{f}=\hat{f}^{+}
\end{equation}
因此，一个实物理量的算符等于它的厄米共轭算符（厄米算符也称为自轭算符）.

对一个厄米算符讲来，属于不同本征值的本征函数彼此正交，现在来讲一下如何直接证明这种正交性.设$f_n$和$f_m$为$\hat{f}$的两个不同本征值，$\varPsi_n$和$\varPsi_m$为其相应的本征函数：
\begin{equation*}
  \hat{f}\varPsi_n=f_n\varPsi_n,\qquad\hat{f}\varPsi_m=f_m\varPsi_m.
\end{equation*}
第一式两边各乘以$\varPsi_m^{*}$，第二式经过复共轭后再在两边各乘以$\varPsi_n$，然后两式相减，得
\begin{equation*}
  \varPsi_m^{*}\hat{f}\varPsi_n-\varPsi_n\hat{f}^{*}\varPsi_m^{*}
  =(f_n-f_m)\varPsi_n\varPsi_m^{*},
\end{equation*}
上式两边对$\mathrm{d}q$积分.由于$\hat{f}^{*}=\tilde{\hat{f}}$，按 (3.14) 上式左边的积分等于零，故得
\begin{equation*}
  (f_n-f_m)\int\varPsi_n\varPsi_m^{*}\,\mathrm{d}q=0,
\end{equation*}
由于$f_n\neq f_m$，得证$\varPsi_n$和$\varPsi_m$两个函数的正交性.

我们在这里只讲了一个物理量$f$，但在本节之初早已说明过，对一组同时可测的物理量的完全集合也可以同样讲.此时完全集合中的每一个量$f,g,\cdots$均有其对应的算符$\hat{f},\hat{g},\cdots$.本征函数$\varPsi_n$则对应于上述诸量同时具有定值的态，即对应于具有一组确定的本征值$f_n,g_n,\cdots$的态，且为下列方程组的共同解：
\begin{equation*}
  \hat{f}\varPsi=f\varPsi,\qquad\hat{g}\varPsi=g\varPsi,\cdots.
\end{equation*}

\section{算符的加法和乘法}

设$\hat{f}$和$\hat{g}$是两个物理量$f$和$g$的算符.它们之和$f+g$对应于一个相应算符$\hat{f}+\hat{g}$.但是，量子力学中不同物理量相加的含义在很大程度上取决于这两个算符能否同时测量.如果$f$和$g$能够同时测量，算符$\hat{f}$和$\hat{g}$就具有共同本征函数，它们也就是算符$\hat{f}+\hat{g}$的本征函数，这一算符的本征值就等于二本征值之和$f_n+g_n$.但当$f$和$g$不能同时具有确定值时，二者之和$f+g$的含义就更有限了.此时我们只能说，算符之和在任一态中的平均值等于二者平均值之和：
\begin{equation}
  \overline{f+g}=\bar{f}+\bar{g}.
\end{equation}
至于算符$\hat{f}+\hat{g}$的本征值和本征函数，一般讲来与$f$和$g$的本征值和本征函数不再有任何关系.如果$\hat{f}$和$\hat{g}$都是自轭算符，显然$\hat{f}+\hat{g}$也是自轭算符，从而它的本征值都是实数，并且等于由此定义的新量$f+g$的值.

下列定理值得提一下.设$f_0$和$g_0$分别为$f$和$g$的最小本征值，而$(f+g)_0$为$f+g$的最小本征值.则可证明
\begin{equation}
  (f+g)_0\geqslant f_0+g_0,
\end{equation}
式中的等号当$f$和$g$可以同时测量时成立.这个定理的证明基于这样一个明显的事实，即一个量的平均值总是大于或等于它的最小本征值.在$f+g$具有本征值$(f+g)_0$的态中，我们有$\overline{f+g}=(f+g)_0$，另一方面，由于$\overline{f+g}=\bar{f}+\bar{g}\geqslant f_0+g_0$，我们就得到不等式 (4.2).

再令$f$和$g$是两个可以同时测量的量.除了它们的和之外，还可以引进它们的乘积概念，这个乘积是这样一个量，这个量的本征值等于$f$和$g$的本征值的乘积.很易证明，这个量的算符作用在函数上的结果，相当于把原先两个量的算符先后作用于该函数上所得的结果.这样的算符，在数学上可表为$\hat{f}$和$\hat{g}$两个算符的乘积.实际上，设$\varPsi_n$为$\hat{f}$和$\hat{g}$的共同本征函数，我们有
\begin{equation*}
  \hat{f}\hat{g}\varPsi_n=\hat{f}(\hat{g}\varPsi_n)
  =\hat{f}g_n\varPsi_n=g_n\hat{f}\varPsi_n=g_nf_n\varPsi_n
\end{equation*}
（记号$\hat{f}\hat{g}$代表一个算符，它作用于函数$\varPsi$上的结果，等于先把算符$\hat{g}$作用于$\varPsi$上再把算符$\hat{f}$作用于函数$\hat{g}\varPsi$后所得的结果）.同样我们也可以用算符$\hat{g}\hat{f}$代替$\hat{f}\hat{g}$，两者的差别仅在于因子的次序.这两个算符作用于$\varPsi_n$函数的结果显然是相同的.由于任一波函数$\varPsi$可表为$\varPsi_n$函数组的线性组合，故$\hat{f}\hat{g}$和$\hat{g}\hat{f}$作用于任一函数的结果也是相同的.这一事实可以用符号等式$\hat{f}\hat{g}=\hat{g}\hat{f}$来表示，或写作
\begin{equation}
  \hat{f}\hat{g}-\hat{g}\hat{f}=0.
\end{equation}
这样的两个算符$\hat{f}$和$\hat{g}$称为可相互对易.由此得出一个重要结论：如果$f$和$g$两个量可以同时取定值，则其算符彼此对易.

上述定理的逆定理也能加以证明（见§11）：如果算符$\hat{f}$和$\hat{g}$对易，那么，它们所有的本征函数都可取成两者的共同本征函数；其物理含义是，这两个算符所对应的物理量可以同时测量.因此算符的可对易性是物理量可以同时测量的一个充要条件.

算符相乘的一个特殊情形是一个算符的幂乘.根据以上的讨论可知，算符$\hat{f}^{p}$（$p$为整数）的本征值等于算符$\hat{f}$的本征值的$p$次方.一般讲来，我们可以把一个算符$\hat{f}$的任意函数$\varphi(\hat{f})$定义为一个算符，这个算符的本征值等于同一函数$\varphi(f)$，其中的$f$是算符$\hat{f}$的本征值.如果函数$\varphi(f)$可展开成泰勒级数，则算符$\varphi(\hat{f})$也能展成这样的幂级数，即归结为各种幂次的$\hat{f}^{p}$.

特别是算符$\hat{f}^{-1}$，它称为$\hat{f}$的逆算符.显然，把算符$\hat{f}$和$\hat{f}^{-1}$先后作用于任一函数上，其结果使该函数保持不变，亦即有$\hat{f}\hat{f}^{-1}=\hat{f}^{-1}\hat{f}=1$.

如果$f$和$g$不能同时测量，它们的乘积概念就不再具有上述直接意义.这一点可用下列事实说明：现在的算符$\hat{f}\hat{g}$不再自轭，从而不能对应于任一实物理量.事实上，根据一个算符的转置定义，我们可写出
\begin{equation*}
  \int\varPsi\hat{f}\hat{g}\varPhi\,\mathrm{d}q
  =\int\varPsi\hat{f}(\hat{g}\varPhi)\,\mathrm{d}q
  =\int(\hat{g}\varPhi)(\tilde{\hat{f}}\varPsi)\,\mathrm{d}q,
\end{equation*}
这里的算符$\tilde{\hat{f}}$只作用在$\varPsi$上而算符$\hat{g}$只作用在$\varPhi$上，所以被积函数不过是$\hat{g}\varPhi$和$\tilde{\hat{f}}\varPsi$两个函数的简单乘积.再一次应用算符的转置定义，我们可写成
\begin{equation*}
  \int\varPsi\hat{f}\hat{g}\varPhi\,\mathrm{d}q
  =\int(\tilde{\hat{f}}\varPsi)(\hat{g}\varPhi)\,\mathrm{d}q
  =\int\varPhi\tilde{\hat{g}}\tilde{\hat{f}}\varPsi\,\mathrm{d}q.
\end{equation*}
于是我们得到一个积分，和原积分相比，其中的函数$\varPsi$和$\varPhi$对调了位置.换句话说，算符$\tilde{\hat{g}}\tilde{\hat{f}}$是$\hat{f}\hat{g}$的转置算符，我们可写成
\begin{equation}
  \widetilde{\hat{f}\hat{g}}=\tilde{\hat{g}}\tilde{\hat{f}},
\end{equation}
即乘积$\hat{f}\hat{g}$的转置，等于其因子分别转置后再以相反次序写出的乘积.(4.4) 式两边都取复共轭，可得
\begin{equation}
  (\hat{f}\hat{g})^{+}=\hat{g}^{+}\hat{f}^{+}.
\end{equation}
如果$\hat{f}$和$\hat{g}$都是厄米算符，则$(\hat{f}\hat{g})^{+}=\hat{g}\hat{f}$.由此可见，当且仅当$\hat{f}$和$\hat{g}$对易时，$\hat{f}\hat{g}$才能是厄米算符.

我们要指出，从两个不对易厄米算符的乘积$\hat{f}\hat{g}$和$\hat{g}\hat{f}$出发，可以对称化成一个厄米算符：
\begin{equation}
  \frac{1}{2}(\hat{f}\hat{g}+\hat{g}\hat{f}),
\end{equation}
这样的表式有时会碰到；称之为对称化乘积.

容易看出$\hat{f}\hat{g}-\hat{g}\hat{f}$是一个反厄米算符（即其转置算符等于其复共轭算符乘上一个负号）.它乘上$\mathrm{i}$后即可变成厄米算符；因此
\begin{equation}
  \mathrm{i}(\hat{f}\hat{g}-\hat{g}\hat{f})
\end{equation}
又是一个厄米算符.

为简便计，今后我们有时将用以下的记号：
\begin{equation}
  [\hat{f},\hat{g}]=\hat{f}\hat{g}-\hat{g}\hat{f},
\end{equation}
并称为这两个算符的对易式.容易证明下列恒等式
\begin{equation}
  [\hat{f}\hat{g},\hat{h}]=[\hat{f},\hat{h}]\hat{g}+\hat{f}[\hat{g},\hat{h}].
\end{equation}
注意，如果$[\hat{f},\hat{h}]=0$和$[\hat{g},\hat{h}]=0$，一般讲并不导致$\hat{f}$和$\hat{g}$对易.

\section{连续谱}

§3和§4中描述离散谱本征函数特性的所有关系式，可以毫无困难地推广到本征值为连续谱的情形中去.

设$f$为具有连续谱的一个物理量.我们可以简单地用同一字母$f$代表它的各个本征值，并用$\varPsi_f$代表相应诸本征函数.按 (3.2) 式，任一波函数$\varPsi$可对离散谱的一套本征函数来展开，与此类似，$\varPsi$也可对具有连续谱的物理量的一套完备本征函数来展开（此时展开式是一个积分式）.这种展开式所具的形状为
\begin{equation}
  \varPsi(q)=\int a_f\varPsi_f(q)\,\mathrm{d}f,
\end{equation}
式中的积分遍及量$f$所能取的整个值域.

连续谱本征函数的归一化问题要比离散谱情形来得复杂.以后将看到，本征函数模量平方积分等于一的要求，现在无法满足.为此，我们把$\varPsi_f$函数的归一化定义改成这样：使得$|a_f|^2\,\mathrm{d}f$等于在$\varPsi_f$波函数所描述的态中测得该物理量的数值介于$f$和$f+\mathrm{d}f$之间的概率.所有$f$可能值的概率总和必须等于1，因此有
\begin{equation}
  \int|a_f|^2\,\mathrm{d}f=1
\end{equation}
[类似于离散谱的 (3.3) 式].

仿效推导 (3.5) 式的方法，应用同一论证，可得下列两式，其一是
\begin{equation*}
  \int\varPsi\varPsi^{*}\,\mathrm{d}q=\int|a_f|^2\,\mathrm{d}f,
\end{equation*}
其二是
\begin{equation*}
  \int\varPsi\varPsi^{*}\,\mathrm{d}q
  =\iint a_f^{*}\varPsi_f^{*}\varPsi\,\mathrm{d}f\,\mathrm{d}q.
\end{equation*}
比较以上两式，发现展开式系数满足下列公式：
\begin{equation}
  a_f=\int\varPsi(q)\varPsi_f^{*}(q)\,\mathrm{d}q,
\end{equation}
此式与 (3.5) 式完全类似.

为了推导归一化条件，现在把 (5.1) 代入 (5.3) 式，结果得
\begin{equation*}
  a_f=\int a_{f'}\left(\int\varPsi_{f'}\varPsi_f^{*}\,\mathrm{d}q\right)\mathrm{d}f'.
\end{equation*}
此式必须对任意的$a_f$都能成立，因而必须是恒等地得到满足.为此，首先，被积函数中的$a_{f'}$的系数（即$\int\varPsi_{f'}\varPsi_f^{*}\,\mathrm{d}q$）只要$f'\neq f$就必须等于零，而当$f'=f$时，这个系数必须变成无穷大（否则对$\mathrm{d}f'$积分后将等于零）.故积分$\int\varPsi_{f'}\varPsi_f^{*}\,\mathrm{d}q$应为$f'-f$的函数，这个函数当宗量不等于零时为零，当宗量等于零时为无穷大，我们用$\delta(f'-f)$代表这个函数：
\begin{equation}
  \int\varPsi_{f'}\varPsi_f^{*}\,\mathrm{d}q=\delta(f'-f).
\end{equation}
我们必须有
\begin{equation*}
  \int\delta(f'-f)a_{f'}\,\mathrm{d}f'=a_f.
\end{equation*}
这个式子确定了$\delta(f'-f)$函数当$f'-f=0$时变成无穷大的方式.由上式显然可得
\begin{equation*}
  \int\delta(f'-f)\,\mathrm{d}f'=1.
\end{equation*}
这样定义的函数称为$\delta$函数（它是由P.A.M.狄拉克首先引入量子力学的）.我们再把定义它的公式写一遍.它们是
\begin{equation}
  x\neq0\ \text{时}\ \delta(x)=0,\qquad\text{而}\ \delta(0)=\infty,
\end{equation}
而且
\begin{equation}
  \int_{-\infty}^{\infty}\delta(x)\,\mathrm{d}x=1,
\end{equation}
式中的上下积分限可换成任意数值，只要$x=0$的点处于积分区内就可以了.设$f(x)$为在$x=0$点处连续的函数，则有
\begin{equation}
  \int_{-\infty}^{\infty}\delta(x)f(x)\,\mathrm{d}x=f(0).
\end{equation}
此式可改写成更普遍的形式
\begin{equation}
  \int\delta(x-a)f(x)\,\mathrm{d}x=f(a),
\end{equation}
式中的积分区间包含$x=a$点，并且$f(x)$在$x=a$点处连续.容易证明$\delta$函数是一个偶函数，即
\begin{equation}
  \delta(-x)=\delta(x),
\end{equation}
最后，根据
\begin{equation*}
  \int_{-\infty}^{\infty}\delta(\alpha x)\,\mathrm{d}x
  =\int_{-\infty}^{\infty}\delta(y)\,\frac{\mathrm{d}y}{|\alpha|}
  =\frac{1}{|\alpha|}
\end{equation*}
可导出
\begin{equation}
  \delta(\alpha x)=\left(\frac{1}{|\alpha|}\right)\delta(x),
\end{equation}
其中$\alpha$为任意常数.

(5.4) 式给出了连续谱本征函数的归一化法则；它可以代替离散谱情形下的归一化条件 (3.6).我们看到，和以前一样，$\varPsi_f$和$\varPsi_{f'}$函数当$f'\neq f$时是相互正交的.可是模量平方$|\varPsi_f|^2$的积分对连续谱讲来是发散的.

函数$\varPsi_f(q)$还满足另一个类似于 (5.4) 式的关系式.要推导这个关系式，可把 (5.3) 式代入 (5.1) 式中，得
\begin{equation*}
  \varPsi(q)=\int\varPsi(q')\left(\int\varPsi_f^{*}(q')\varPsi_f(q)\,\mathrm{d}f\right)\mathrm{d}q',
\end{equation*}
由此我们可以立刻推知必有
\begin{equation}
  \int\varPsi_f^{*}(q')\varPsi_f(q)\,\mathrm{d}f=\delta(q'-q).
\end{equation}
离散谱情形中当然也有一个与上式相类似的关系式：
\begin{equation}
  \sum_n\varPsi_n^{*}(q')\varPsi_n(q)=\delta(q'-q).
\end{equation}
(5.1)，(5.4) 两式和 (5.3)，(5.11) 两式相比较，可以看出，一方面，函数$\varPsi(q)$可按函数组$\varPsi_f(q)$展开，其展开系数为$a_f$，另一方面，(5.3) 式是一个完全类似的展开式，式中的$a_f\equiv a(f)$函数按函数组$\varPsi_f^{*}(q)$展开，而$\varPsi(q)$充当了展开式系数.函数$a(f)$和$\varPsi(q)$一样，也能完全确定此系统的状态；我们有时把$a(f)$称为$f$表象中的波函数[而$\varPsi(q)$称为$q$表象中的波函数].正如$|\varPsi(q)|^2$确定系统的坐标值介于给定的$\mathrm{d}q$区间内的概率一样，$|a(f)|^2$确定了$f$的数值介于给定的$\mathrm{d}f$区间内的概率.一方面，$\varPsi_f(q)$函数为物理量$f$在$q$表象中的本征函数；另一方面，它的共轭复函数$\varPsi_f^{*}(q)$就是坐标$q$在$f$表象中的本征函数.

设$\varphi(f)$为物理量$f$的某种函数，并且$\varphi$和$f$的关系是一一对应的.那么，每个$\varPsi_f(q)$函数都可以看作$\varphi$的一个本征函数.可是，这时这些函数的归一化问题必须改变，$\varphi$的本征函数$\varPsi_{\varphi}(q)$应按以下条件归一化：
\begin{equation*}
  \int\varPsi_{\varphi(f')}\varPsi_{\varphi(f)}^{*}\,\mathrm{d}q
  =\delta[\varphi(f')-\varphi(f)].
\end{equation*}
而$\varPsi_{f'}$函数是按条件 (5.4) 归一化的.$\delta$函数只有当宗量$f'=f$时才等于零.当$f'$趋近$f$时，我们有
\begin{equation*}
  \varphi(f')-\varphi(f)=\frac{\mathrm{d}\varphi(f)}{\mathrm{d}f}(f'-f).
\end{equation*}
按 (5.10) 式我们可写成\footnote{一般来讲，如$\varphi(x)$为某一单值函数（其逆函数无需单值），我们就有
\begin{equation*}
  \delta[\varphi(x)]=\sum_i\frac{1}{|\varphi'(\alpha_i)|}\,\delta(x-\alpha_i),
\end{equation*}
其中$\alpha_i$为方程式$\varphi(x)=0$的各个根.}
\begin{equation}
  \delta[\varphi(f')-\varphi(f)]
  =\frac{1}{\left|\dfrac{\mathrm{d}\varphi(f)}{\mathrm{d}f}\right|}\,
  \delta(f'-f).
\end{equation}
此式和 (5.4) 式对比，可知$\varPsi_{\varphi}$和$\varPsi_f$函数的相互关系为
\begin{equation}
  \varPsi_{\varphi(f)}
  =\frac{1}{\sqrt{\left|\dfrac{\mathrm{d}\varphi(f)}{\mathrm{d}f}\right|}}\,
  \varPsi_f.
\end{equation}
还有一些物理量，这种量在一个值域内具有离散谱，而在另一值域内具有连续谱.对这种物理量的本征函数讲来，本节和前节中导出的所有公式当然也能成立，但有一点必须指出，它的完备函数组是由离散谱和连续谱的本征函数加在一起组成的.因此，任一波函数对这种量的本征函数组展开时具有以下的形式：
\begin{equation}
  \varPsi(q)=\sum_n a_n\varPsi_n(q)+\int a_f\varPsi_f(q)\,\mathrm{d}f.
\end{equation}
式中对离散谱求和，并对整个连续谱求积分.

坐标$q$本身是具有连续谱的物理量的一个例子.容易证明：它所对应的算符相当于简单地乘以因子$q$.由于各种不同坐标值的概率是由模量平方$|\varPsi(q)|^2$确定的，所以坐标平均值应为
\begin{equation*}
  \bar{q}=\int q|\varPsi|^2\,\mathrm{d}q
  \equiv\int\varPsi^{*}q\varPsi\,\mathrm{d}q.
\end{equation*}
将上式和算符定义 (3.8) 式比较，得到\footnote{为简单计.今后我们常把恒等于乘以某因子的算符直接写成该因子.}
\begin{equation}
  \hat{q}=q,
\end{equation}
根据一般规则，这个算符的本征函数应该由方程$q\varPsi_{q_0}=q_0\varPsi_{q_0}$确定，式中的$q_0$暂时代表一个具体的坐标值，以便和变量$q$相区别.由于$\varPsi_{q_0}=0$或$q=q_0$时这个方程式都能得到满足，很明显，满足归一化条件的本征函数应该是\footnote{任一$\varPsi$函数对这种本征函数展开后的展开系数为
\begin{equation*}
  a_{q_0}=\int\varPsi(q)\delta(q-q_0)\,\mathrm{d}q=\varPsi(q_0)
\end{equation*}
故坐标值介于给定区间$\mathrm{d}q_0$内的概率为
\begin{equation*}
  |a_{q_0}|^{2}\,\mathrm{d}q_0=|\varPsi(q_0)|^{2}\,\mathrm{d}q_0
\end{equation*}
这正是应有的结果.}
\begin{equation}
  \varPsi_{q_0}=\delta(q-q_0).
\end{equation}

\section{过渡到经典力学极限情形}

量子力学把经典力学作为自己的某种极限形式包括在内.问题是如何过渡到这种极限情形.

在量子力学中，一个电子是由波函数描述的，波函数确定了电子坐标的各种数值；对于这种波函数，迄今我们只知道它是某种线性偏微分方程的解.另一方面，经典力学中的一个电子被看成一个质点，运动于由运动方程所完全确定的轨道上.量子力学和经典力学之间的这种内在关系，在某种意义上，与电动力学中波动光学和几何光学之间的内在关系相类似.在波动光学中，电磁波是由满足一定的线性微分方程组（即麦克斯韦方程）的电场矢量和磁场矢量所描述的.但在几何光学中，光的传播被看作是沿着确定的轨道（光线）行进的.这种类似性，使我们能像波动光学过渡到几何光学那样，将量子力学过渡到经典力学的极限情形.

我们来回忆一下波动光学在数学形式上是如何过渡到几何光学的（见《场论》，§53）.设$u$为电磁波的任一场分量.它可表成$u=ae^{i\varphi}$的形式（$a$和$\varphi$为实量），其中的$a$称为波的振幅，$\varphi$称为波的相位.几何光学对应于波长很短时的极限情形；从数学上讲来，就是相位$\varphi$（在几何光学中，$\varphi$称为程函）在短距离内具有很大的改变量；这就是说，此时可以假定$\varphi$本身具有很大的绝对值.

与此类似，我们可以从这样的假定出发，假定量子力学的波函数在经典力学极限情形下具有$\varPsi=ae^{i\varphi}$的形式，式中的$a$是一个缓变函数，而$\varphi$取很大的数值.大家知道，力学中的质点轨道可以通过变分原理确定，按此原理，一个力学系统的作用量$S$必须取极小值（最小作用量原理）.在几何光学中，光线是由所谓费马原理确定的，按此原理，光线的光程亦即轨道的起端与终端的相位差必须取最小（或最大）可能值.

基于这一类似性，我们可以断言，在经典极限情形下，波函数中的相位$\varphi$应与所考虑物理系统的力学作用量$S$成正比，即有$S=\text{常量}\times\varphi$.这个比例常量称为普朗克常量，我们用$\hbar$表示\footnote{1900年由M.普朗克引入物理学中.本书各处所用的$\hbar$，严格来讲应等于普朗克常量$h$除以$2\pi$；$\hbar$是狄拉克最先使用的.}.$\hbar$具有作用量的量纲（因$\varphi$量纲为1），并且
\begin{equation*}
  \hbar=1.055\times10^{-34}\ \mathrm{J}\cdot\mathrm{s}
\end{equation*}
因此一个“几乎经典的”（或称准经典的）物理系统的波函数具有下列形式：
\begin{equation}
  \varPsi=ae^{iS/\hbar}.
\end{equation}
普朗克常量$\hbar$在一切量子现象中占有重要地位.它的相对值（与同量纲的其它物理量相比）决定着该物理系统的“量子化程度”.量子力学向经典力学的过渡，相当于相位很大时的情形，这可用$\hbar$趋于零（$\hbar\to0$）的形式描述（正如波动光学过渡到几何光学时，相当于波长趋于零，$\lambda\to0$）.

我们已经阐明了波函数的极限形式，但未涉及它与经典轨道运动的关系问题.一般讲来，用波函数描述的运动并不趋向确定的轨道运动.它与经典运动的关系是这样的，假定在某一起始时刻，波函数和随之而有的坐标概率分布值已给定，在随后各时刻，这个概率分布将按经典力学的定律变化（详见§17的末段）.

为了得出沿确定轨道的运动，我们必须从特殊形状的波函数出发，这种波函数只在一个很小的空间范围内才显著地不等于零（称为波包）；这个空间范围的尺度必须随$\hbar$一起趋于零.然后我们才能说，在准经典情形下，该波包在空间将沿质点的经典轨道运动.

最后，量子力学算符在这种极限情形下应该还原成为一个相乘因子，这个因子就是相应的物理量.

\section{波函数与测量}

让我们再回到测量过程，它的性质已经在§1中定性地讨论过；现在来说明这些性质和量子力学的数学表述是怎样联系的.

考虑包括以下两个部分的一个系统：一个经典仪器和一个电子（看成一个量子客体）.测量过程就是这两个部分进入相互作用，其结果是仪器由初态转为另一状态，根据其状态变化可以引出电子状态的有关结论.仪器的状态是由标志它的某种（或某些）物理量的数值——即“仪器的读数”——所确定的.我们暂时用$g$表示这个量，用$g_n$表示$g$的本征值；$g$的值域按照仪器的经典性质一般讲来应该是连续的，但我们——仅为今后简化公式起见——假定它为离散谱.仪器的状态可用准经典波函数$\varPhi_n(\xi)$描写，下标$n$对应于仪器“读数”$g_n$，而$\xi$为其坐标的集合.仪器的经典性质表现于下列事实，在任一给定时刻，我们可以肯定它处于具有确定$g$值的某一已知态$\varPhi_n$中；这样的假定，对一个量子系统讲来，当然是不合理的.

令$\varPhi_0(\xi)$为仪器的初态波函数（测量之前），$\varPsi(q)$为电子的任一归一化初态波函数（$q$表示其坐标的集合）.这两个函数相互独立地描述仪器和电子的状态.故整个系统的初态波函数等于下列乘积：
\begin{equation}
  \varPsi(q)\varPhi_0(\xi),
\end{equation}
随后，仪器和电子进入相互作用.运用量子力学的方程式，可在原则上追踪该系统的波函数随时间的变化.在测量过程结束之后，这个波函数当然不一定仍是$\xi$的函数和$q$的函数的乘积.把这个波函数对仪器的本征函数$\varPhi_n$（它们构成一完备组）展开后，可得下列叠加形式：
\begin{equation}
  \sum_n A_n(q)\varPhi_n(\xi),
\end{equation}
其中$A_n(q)$为$q$的某种函数.

仪器的“经典性质”，以及经典力学既是量子力学的基础又是其极限情形的双重角色，现在开始露面.如前所述，仪器的经典性质意味着，$g$量（“仪器读数”）在任何时刻均有某种定值.这就使我们能够断定，测量之后，仪器加电子的整个系统的态，实际上并不是由 (7.2) 式的整个级数之和描述，而是由其中与仪器“读数”$g_n$相对应的那一项描述的：
\begin{equation}
  A_n(q)\varPhi_n(\xi),
\end{equation}
由此可见，上式中的$A_n(q)$正比于测量结束后电子的波函数.从函数$A_n(q)$并没有归一化这一点也可以看出，它还不是电子波函数本身，$A_n(q)$中不但包含有关电子末态性质的信息，而且还包含着仪器出现第$n$个“读数”的概率（由系统初态所确定）的信息.

由于量子力学方程的线性性质，$A_n(q)$和电子的初态波函数$\varPsi(q)$的关系，一般讲来，可以通过某种线性积分算符表出：
\begin{equation}
  A_n(q)=\int K_n(q,q')\varPsi(q')\,\mathrm{d}q',
\end{equation}
其中的核$K_n(q,q')$表述这个测量过程的特征.

我们假定上述测量给出了电子态的完全描述.换句话说（见§1），末态中所有各量的概率应该与电子的先前（测量以前）状态无关.从数学上讲，这意味着函数$A_n(q)$的形状必须由测量过程本身所确定，而与电子的初态波函数$\varPsi(q)$无关.因此$A_n$应呈下列形式：
\begin{equation}
  A_n(q)=a_n\varphi_n(q),
\end{equation}
其中$\varphi_n$为某种确定的函数，我们假定它已经归一化，只有常数$a_n$才依赖于$\varPsi(q)$.在积分关系式 (7.4) 中，这相当于核$K_n(q,q')$分解成为$q$的函数和$q'$的函数二者的乘积
\begin{equation}
  K_n(q,q')=\varphi_n(q)\varPsi_n^{*}(q').
\end{equation}
从而常数$a_n$和函数$\varPsi(q)$的线性关系呈下列形式：
\begin{equation}
  a_n=\int\varPsi(q)\varPsi_n^{*}(q)\,\mathrm{d}q,
\end{equation}
其中$\varPsi_n(q)$为依赖于测量过程的某些确定的函数.

函数$\varphi_n(q)$就是测量之后电子的归一化波函数.这里我们看到，用测量方法确定一个电子态（这个电子态被一个确定的波函数描述）的可能性在理论的数学表述中是怎样得到反映的.

如果对一个波函数给定为$\varPsi(q)$的电子进行这种测量，则常数$a_n$具有一个简单的物理含义：根据一般规则，$|a_n|^2$就是测得第$n$个结果的概率.所有测量结果的总概率应该等于一：
\begin{equation}
  \sum_n|a_n|^2=1.
\end{equation}
为了使 (7.7) 式和 (7.8) 式对任意的归一化的函数$\varPsi(q)$都能成立，任意函数$\varPsi(q)$必须能对函数组$\varPsi_n(q)$展开（参考§3），这就是说，$\varPsi_n(q)$构成一套正交归一化的完备函数组.

如果电子的初态波函数等于其中的一个函数$\varPsi_n(q)$，则与$\varPsi_n$对应的常数$a_n$显然等于1，而其它的$a_n$都等于零.换句话说，对处于$\varPsi_n(q)$态的电子进行该种测量后，肯定能够得出第$n$个结果.

函数$\varPsi_n(q)$的上述一切性质表明，它们就是用以标志电子的某种物理量（用$f$表示）的本征函数，我们所讲的那种测量，可以说是对$f$这个量进行测量.

十分重要的是，$\varPsi_n(q)$这组函数一般讲来并不等于$\varphi_n(q)$这组函数.一般讲来，后一组函数甚至并不相互正交，也不构成任一算符的本征函数组.这就表明了这样一个事实，量子力学中的测量结果是无法重现的.如果电子处于$\varPsi_n(q)$态，那么，测量$f$的结果，会肯定地得出$f_n$值.但当测量结束之后，该电子就处于不同于初态的$\varphi_n(q)$态，一般讲来，在这个态中$f$不再取任何定值.故在第一次测量结束之后，对该电子紧接着作第二次测量时，我们会得到不同于第一次测得的$f$值\footnote{关于测量的无法重现性问题，必须指出一个重要例外——测量结果可以重现的一个量就是坐标.在足够短的时间间隔内对一个电子的坐标进行两次测量，一定会得出相邻值；否则将意味着电子具有无限大的速度.从数学上讲来，这一点与下列事实有关，即电子和仪器的相互作用能量算符与坐标算符互相对易，因为这种相互作用能（在非相对论理论中）仅为坐标的函数.}.要想用第一次测量的已知结果去预断第二次测量的结果（意即算出其概率），我们必须首先知道：由第一次测量所建立的那个态的波函数$\varphi_n(q)$，以及在第二次测量中欲求其概率的那个波函数$\varPsi_n(q)$.这就是说，从第一次测量结束时的$\varphi_n(q)$出发，应用量子力学方程求出$\varphi_n(q,t)$波函数；那么，在$t$时刻进行第二次测量时获得第$m$个结果的概率可以用积分$\int\varphi_n(q,t)\varPsi_m^{*}(q)\,\mathrm{d}q$的模量平方给出.

我们看到，量子力学中的测量过程具有“两面”性：它对电子的过去和未来起着不同的作用.对过去而言，即对前次测量中所建立的一个电子态而言，通过这次测量，可以“验证”由该态所预断的各种可能结果的概率.对未来而言，通过这一次的测量后又建立了一个新的电子态（参见§44\mbox{）.}因此，测量过程本身的这个特性包含着一个深邃的不可逆性原理.

这种不可逆性具有重要的原则意义.我们将在以后看到（见§18末段），量子力学基本方程本身对时间的变号具有对称性；从这一方面讲来，量子力学和经典力学没有什么区别.但是测量过程的不可逆性，使得时间的两个指向具有物理上的不等价性，也就是说，使得过去和未来呈现差别.
